% Copyright (c) Chris Paxton 2026
% Licensed under the Apache License, Version 2.0.
\section{Whole-body trajectory control: interface and qualification protocol}
\label{app:whole_body_control}

\paragraph{Task-space contract.}
The control interface consists of a masked base pose in $SE(2)$, one masked
end-effector pose in $SE(3)$ per hand, and a vector of gripper coordinates per
hand. Mars and Stretch are the initial single-hand implementations. Inactive
base constraints permit coordinated motion to satisfy an active hand target.
The policy observes proprioceptive and command history together with a timed
task-space preview; it does not observe simulator contact forces, object poses,
or the joint-space witness used to construct reachable targets. The executor
owns grasp/release transitions. Direct-velocity policies command bounded base
and arm velocities; demonstration initialization remains optional rather than
requiring a residual policy architecture.

\paragraph{Related work and research direction.}
Blind Precise Assembly~\cite{bronars2026blindpreciseassemblyproprioceptive}
learns local assembly from proprioceptive history using identified actuator
dynamics, training-only contact supervision, and a privileged critic. Real-world
peg insertion is demonstrated on a seven-DoF arm; gear and nut tasks are shown
in simulation. Its contact-starting local skill does not establish Mars grasp
retention, mobile-base safety, or damage prevention. Its mjfit generator preserves
assembly clearances in collision geometry rather than identifying robot dynamics.
These motivate future observability and isolated training ablations, while we
retain TAMP goal conditioning and hard safety limits; no such ablation is claimed
as implemented or validated here.

\paragraph{References and physical screening.}
Ground-truth TAMP planning supplies collision-checked arm paths, retimed with
minimum-jerk segments under robot-profile speed, acceleration, and jerk limits.
Offline FK constructs reachable hand targets; no joint-reference hint is
provided to the policy. Mars setup candidates preserve the world grasp pose
while considering different base/arm configurations. Loaded static support
margin is screened along the lift and carry path. This screening is not a
certificate of dynamic feasibility: execution must independently satisfy
retention, tracking, contact-load, and wheel-support checks. Payload attachment
is restricted to scratch planning state; continuous evaluation uses physical
gripper contacts without attachments or stage resets.

\paragraph{Admission and execution are separate outcomes.}
An initial ten-case Mars carry diagnostic admitted only two references; both
executed successfully. Replays identified support-margin failures already at
the initial path sample, alongside an inconsistent reset lift height.
This motivated explicit rejection reasons, path-aware setup selection, and
consistent configured lift heights. Subsequent development-selected checks
admitted and executed ten of ten carry references. A separate fifteen-case
smoke test covered 20, 50, and 100\,g box payloads across translation,
orientation, combined motion, endpoint changes, and small windowed disturbances;
all fifteen passed with the feedback controller. These changed development
distributions are not a paired improvement estimate, a learned-policy result,
or a general reliability estimate. Continuous pick/place regressions remain
a separate gate; carry-component success does not establish that gate.

\paragraph{Source-bound teacher qualification.}
At revision \texttt{0101dc35}, all 300 carry references were admitted and
297 executed successfully: 100/100 at 20\,g, 99/100 at 50\,g, and 98/100 at
100\,g. One 50\,g disturbance case tripped the wheel-unloading guard; two
100\,g wrist-roll cases exceeded the unchanged 0.1\,rad relative-object rotation
limit. A separate continuous pick/return-place suite passed 30/30 cases on one
pedestal. The wheel event blocked training despite the high average success.
These feedback results and failures are retained under
\texttt{paper/data/whole\_body\_control/qualification\_v2.json}.
The subsequent setup revision chooses maximum support margin for disturbed
carries; the exact failing disturbance passed a targeted replay without changing
its wrench. Full requalification is required before starting PPO. Grasp-depth,
contact-impedance, and timing ablations did not reliably resolve the roll-slip
failures; no relaxed retention threshold or increased torque was adopted.
Full requalification at \texttt{ba658820} then admitted 300/300 references and
completed 298/300 carries (100/100, 100/100, and 98/100 by increasing mass), with
zero unsafe-contact or wheel-unloading terminations. Continuous pick/return-place
again passed 30/30. The two 100\,g roll-slip failures remained counted. These
results met the then-current simulator teacher gate and are retained in
\texttt{paper/data/whole\_body\_control/qualification\_v3.json}; they do not
qualify the subsequently trained policy or demonstrate hardware compliance. Subsequent
review tightened source, catalog and regression-test binding; these historical
reports are retained as observations, not authorization for training under the
revised sources. The reported gate is per mass and conditional on admission,
not a 95\% per-stratum or end-to-end success guarantee.

\paragraph{Stage continuity and command-interface checks.}
An explicit stage-rearming refactor retained the 298/300 carry result but reduced
continuous pick/return-place success to 27/30: one retention loss and two tracking
timeouts, all during placement. Development-selected ablations isolated loss of
executor-authorized base compensation across the transition; restoring tracking
trim alone did not recover all three cases. Explicitly authorizing compensation
for verified free-space descent, while disabling it at support contact and on
faults, restored 30/30 cases with unchanged contracts and zero unsafe-contact or
wheel events. The subsequent immutable \texttt{b181db61} rerun completed 298/300
carries and 30/30 continuous cases, with 367 fresh source-bound regression tests;
the combined canonical readiness check passed for that snapshot. Both remaining
100 g orientation retention failures are preserved. Complete evidence is archived
in \texttt{qualification\_v6.json} and adjacent compressed reports; it does not
qualify later public-tool or controller changes. A simulator command-protocol prototype also replayed four
admitted TAMP arm paths, including loaded carry, using task-pose schedules and
a measured-Jacobian teacher. This component result does not establish full
public-tool, learned-policy, or hardware acceptance. Complete placement and
component evidence is archived as \texttt{placement\_recovery\_v6} and
\texttt{task\_pose\_command\_smoke} in the control data directory.

The public Mars adapter now connects task discovery, private collision-checked
planning and independently stepped physical execution on two distinct supports.
A separate frozen thirty-case catalog spans two layouts, five seeds and three
payload masses. All cases have reachable fixture witnesses. A selected 20 g
transfer completed with 11.8 cm measured base displacement, 0.013 rad peak relative
object rotation, 3.24 N peak contact load and stable support after release, without
wheel-lift termination. Its roughly 100 s simulated execution used explicit
geometric-path retiming and a bounded support-contact approach; it is not a
fixed-deadline or full-catalog result. Heavier-payload retention remains open. A
separate harness uses the normal agent prompt, parser and dispatcher, gated on
oracle success and normal observed inputs. No actual observed-agent trial or
hardware outcome is claimed from these interface tests.

The subsequent complete public-tool run at \texttt{9dee06fc} improved from 4/30
to 16/30 on that unchanged catalog: 10/10 at 20 g, 6/10 at 50 g, and 0/10 at
100 g, with all thirty fixtures admitted and 588 source-bound regression tests
passing. The remaining fourteen outcomes comprise six retention losses, five
wheel-unloading stops, one base-drift stop and two planning failures. Successful
tasks took a median 67.55 simulated seconds. Shorter base approaches and verified
steady holds improved the baseline; heavier grasp/transport reliability remains
unqualified. Full traces, failures and identities are archived in
\texttt{public\_tamp\_v3.json} and its adjacent compressed report. These frozen
benchmark results are not an untouched held-out generalization study.

\paragraph{Qualification and learning.}
The held-out protocol fixes 100 cases per payload mass, balanced across the five
motion categories. It reports admission, success conditional on admission,
end-to-end success, pose errors, settling time, post-completion drift, retention,
jerk, contact loads, and wheel unloading. Training readiness requires at least
95\% admission and 95\% conditional feedback execution success per mass, no
unsafe-contact or wheel-lift terminations, and success on a separate thirty-case
continuous pick/place regression. Learned checkpoints are compared with the
feedback controller on identical cases. Progress and tracking rewards are
combined with smoothness, effort, contact, overshoot, and settled-motion costs;
randomization increases only after validation. Hardware force calibration,
full public-tool Mars task execution, and qualification of larger native
objects remain outside the evidence reported here.

An optional simulator research mixture reserves 70\% of resets for ordinary
tracking, with 10\% pregrasp, 12\% carry, 5\% placement and 3\% disturbance.
Loaded resets vary payload identity, mass, center of mass and object friction;
closed-empty tracking trials distinguish a retention command from grasp success.
A versioned observation suffix provides retention intent and gripper velocity.
An advisory object-in-hand estimator learns from debounced bilateral-contact
labels using encoder and command history, with a separate supervised optimizer
and episode-disjoint audit samples. Its loss does not alter the PPO reward or
share actor gradients. Contact occupancy does not certify force closure or
future retention; probability calibration and no-regression task qualification
remain required before using this estimate in execution decisions.

Development composition tests preserve simulation state between transfers and
compare scripted selection with a real model using the same public TAMP tools.
A selected Stretch two-object regional task passes for both selectors. A Mars
delivery-and-return task completes the first transfer but exhausts the return
planner's budget for both selectors, exposing a post-placement composition gap
that isolated tracking misses. These privileged, independently stepped pilots
use reference controllers.

A subsequent simulator execution backend connects a fixed learned checkpoint
to these public tools without teacher fallback. Training-source/weight binding,
reconstruction of the original training-model fingerprint, and a separate
compiled robot/actuator/mesh/solver comparison admit an additional fixed support
without disabling model checks. Three matched development cases span two
layouts and 20/50\,g payloads. The reference completes delivery in all three,
then fails return planning; the checkpoint completes zero deliveries, making
14 actor calls per case before the unchanged fixed-base arm certificate stops
approach drift. All three stops are measured and confirmed, with no wheel-lift
or force-limit termination. These failures establish an exercised interface,
not useful learned pick/place performance or compliant contact. Grounding,
live asynchronous execution, broader household tasks and hardware remain open.

A subsequent startup audit finds that these initial failures preceded the first
arm command. An explicit idle actuator hold fixes that lifecycle error; the
same checkpoint then reaches a commanded pregrasp target but still stops on
base drift in all three cases, after 13/14/13 actor calls. In a selected
stationary fixture, reference and zero-velocity controls keep hand position
error below 0.5\,mm, while the checkpoint reaches approximately 15\,mm even
after one second of zero-velocity history. Inherited carry-dominated statistics,
unseen task-phase categories and missing stationary-base manipulation sampling
motivate further training ablations. Geometric replanning still uses IK during
public execution; policy inference itself does not.

\paragraph{Hold-focused policy ablation.}
Subsequent simulator pilots add stationary-base manipulation resets, physical
pregrasp holds, categorical identity normalization and a six-DOF body-motion
cost while retaining 70\% ordinary tracking resets. An optional learned
feedback head computes bounded mean motion from a context-dependent gain matrix
applied to current/future pose errors and measured velocities. Consequently its
deterministic motion mean is zero when these signals are zero; this does not
prove stability or convergence. A separate bounded gripper head retains release
authority constraints. Tracking context/critic transfer explicitly reinitializes
this head and the PPO optimizer. Five bounded runs add 24,576 PPO samples.

The selected 12,288-sample feedback candidate holds hand position within
1.17\,mm over ten seconds on five development pregrasp fixtures, with peak base
displacement 1.70\,mm. None satisfies a continuous two-second six-DOF quiet hold.
Nearby target residuals are 6.40--7.48\,mm after three seconds and
3.81--6.08\,mm after six seconds; reachable 0.12\,rad wrist targets retain
0.118--0.122\,rad error. One rotation trial stops on an inference deadline.
The same candidate makes 101 actor calls on each of three public TAMP programs
but completes no transfer: the first joint certificate rejects
0.0153--0.0158\,rad residual despite 1.39--1.45\,mm hand error. This exposes
an outstanding Cartesian-policy/joint-executor contract mismatch. The broader
paired no-regression gate also fails. Raw traces, source/weight bindings and
matched teacher results accompany \texttt{hold\_feedback\_*\_v2.json}; a fresh
692-test regression receipt covers the implementation. No checkpoint or
hardware capability is qualified by these development measurements.

\paragraph{Feedback and execution-contract diagnosis.}
Matched frozen pregrasp ablations identify an additional wheel velocity loop as
a source of stationary chatter: reducing its gain from 2.0 to 0.5 produces
5.85--6.00\,s continuous quiet in two six-second, fixed-weight trials, without
changing robot geometry or solver parameters. A refined actor subsequently
passes the two-second six-DOF quiet criterion in five stationary development
fixtures. These results do not establish moving-goal or loaded-task safety.
New training options align fixed-base episodes with the executor's
1\,cm / 0.02\,rad drift boundary and add dense hand progress evaluated against
the same active target at both ends of a transition, alongside fine pose,
jerk, force and retention terms.

An explicit simulator Cartesian arm mode consumes the complete FK schedule
from a planned joint witness. It checks measured configurations, static support,
retention and base drift each control cycle, then tests the original requested
Cartesian endpoint rather than an arbitrary IK joint solution. A bounded final
pose correction avoids adding IK and execution tolerances. These sampled checks
are not a continuous collision certificate for the learned configuration path.
The first three public development programs still fail pregrasp tracking in
this mode, making 371/391/371 actor calls; the joint mode fails after 101 each.
Later refinements and teacher comparisons are preserved with all failures in
\texttt{hold\_cartesian\_validation\_v4.json}. No learned pick/place or hardware
qualification follows from these interface and stationary-hold improvements.

An additional 8,192-sample refinement yields 3.20/3.47\,mm nearby endpoint errors
on two development fixtures, with 0.0455/0.0479\,rad angular residuals. Neither
meets the stricter public arm contract. Learned public execution remains 0/3.
An explicit analytic Cartesian teacher completes pregrasp, grasp and lift in a
frozen 20\,g case, then exposes a terminal hold-transition unloading transient.
Preserving the first safety cause avoids misreporting this as a timeout.
Clearing transport compensation and waiting for measured quiet before position
hold removes the detected unloading in this selected replay. Both tested
analytic rotation metrics complete delivery with independently verified support,
release and stability; both then exhaust the return planner budget. Maximum
wheel gaps are 0.031--0.038\,mm under the unchanged 0.10\,mm unloading threshold,
with peak reported contact force approximately 3.73\,N. These are repeated
analytic controls on one physical case, not learned or population qualification.
All original and repaired outcomes are retained in
\texttt{cartesian\_pose\_oracle\_v5.json}.


\paragraph{Precision refinement and reference progression.}
A further 65,536 PPO samples reduce nearby position residuals to
0.420/0.402\,mm on two development cases, while angular residuals of
0.0245/0.0232\,rad still miss the 0.02\,rad contract. Two wrist probes avoid
the earlier base-envelope abort but retain 0.0585/0.0539\,rad angular error.
Actor-input auditing also identifies and repairs a Cartesian stage handoff that
incorrectly retained the release phase. A binary reference-clock gate freezes
at approximately 3.1\,mm intermediate lag; continuous lag scaling restores
progress within the unchanged 10\,mm / 0.10\,rad tracking tube. The final
3\,mm / 0.02\,rad endpoint and all physical limits remain unchanged.
After both repairs, learned execution still completes no transfers in three
public programs: two pregrasp timeouts and one tracking-envelope stop, all
confirmed. A fresh analytic comparison completes delivery with the 0.15\,m/rad
metric but hits a cycle deadline with 1.0\,m/rad; both outcomes are retained.
These results separate integration defects from remaining learned tracking and
runtime limitations. Full traces, negative paired comparisons and original
training sources accompany \texttt{pose\_precision\_training\_v6},
\texttt{pose\_precision\_validation\_v6} and \texttt{cartesian\_clock\_v7}.


\paragraph{Continuation regression and broader validation.}
Two further 16,384-sample pilots compare ordinary and goal-gated settling.
Both retain stationary quiet but regress on moving goals relative to their
initial checkpoint. Goal gating improves nearby error relative to ordinary
continuation, yet all six public replays stop during pregrasp at the unchanged
tracking bound. On identical recorded observations, a weights/normalization
counterfactual audit finds larger command changes from updated weights than
normalization alone; it does not establish counterfactual rollout success.
The optimizer was restarted at five times the source run's final learning rate.
A subsequent factorial diagnostic separates restart rate and fixed/adaptive
normalization before changing the task distribution. Its validation includes
ordinary base, hand and combined goals, retargets and bounded pushes alongside
TAMP. Current retention-error running cost is corrected independently of
peak-based safety evidence. Neither accurate stationary holds nor additional
training alone establishes task qualification.


\paragraph{Recovered continuation and reachable translations.}
All four learning-rate/normalization pilots complete before a host reboot;
checksummed numbered checkpoints and full negative evaluations survive.
Fixed statistics with smaller updates preserve nearby position accuracy
(0.29--0.39\,mm), but angular residuals remain 0.026--0.028\,rad and public
execution remains unsuccessful. Correcting base-first action grouping in an
offline audit leaves measured outcomes unchanged and attributes larger command
changes to weights than normalization. A bounded additional curriculum samples
angular-nullspace directions and verifies finite FK witnesses before ordinary
collision/derivative admission. It adds 5--20\,mm translations with unchanged
hand attitude, rather than relying only on random joint goals. On eight new
translation development cases, the teacher achieves hand accuracy and settles
at least once in every case but retains final settling in five. The original
learned actor succeeds in one and trips the unchanged unloading guard in one.
These are diagnostic gaps, not qualification or infeasibility labels. General
base/hand/combined, retarget and push probes remain alongside TAMP in the next
matched comparison; no runtime IK or safety-limit relaxation is introduced.

\paragraph{Integrated agent delivery diagnostic.}
The matched translation pilots subsequently complete 16,384 PPO steps each
without improving public task success (0/3 programs per actor). A separate
single-delivery development contract exercises the actual remote agent through
scene discovery, planning and execution tools. Reference execution completes
the delivery and independent quiet verification; the learned backend does not.
Explicit offline arm-feedback scaling adds no PPO steps and allows two variants
to reach grasp, exposing a lift-start payload/support collision that had been
misreported as IK failure. Per-seed diagnostics now preserve successful IK
followed by path rejection. An opt-in Cartesian lift rule permits separation
from existing shallow static support contact (at most 2\,mm initial compression,
no deepening); other collision and measured safety checks remain active.
With this distinct contact contract, the two variants begin lifting but stop on
orientation-tube and base-drift limits respectively. No transfer completes.
These development failures motivate loaded-transition training and do not
qualify either the controller or the new contact behavior for hardware.
The final-source comparison passes 3/3 reference deliveries and 0/3 learned
deliveries from matched initial physics; the first fixture uses actual agent
selection for both backends, and the other two use scripted public-tool
selection. All learned failures occur during lift and confirm a measured stop.

\paragraph{Supported grasp-to-lift coverage.}
Ordinary carry resets finish a teacher lift before policy handoff, while
pregrasp episodes reach lift only after a successful learned approach. A new
opt-in subtype therefore starts at physical supported closure and supplies a
complete smooth lift/hold pose schedule. The proposed mixture assigns it 9\%
of resets within the existing carry share, retaining 70\% tracking and the
remaining manipulation tasks. Quiet closure, private path/support admission and
measured collision checks precede and accompany rollout. Cylindrical payloads
still slide/spin during teacher load transfer, including after closure and depth
diagnostics; their failures remain open. The initial subtype therefore uses
boxes, while other tasks retain their object mix. Eight box preflight cases
admit, seven settle at least once, and five retain final success, with one
retention loss and two settling failures. A 512-step PPO smoke test completes;
these are runtime/distribution checks, not controller qualification or evidence
that the subsequent matched training improves public task success.
